\chapter{Predicción de la demanda de transporte público: dominio, técnicas y trabajos previos}\label{cap:estado_arte}

Tres bloques temáticos trazan la línea argumental que justifica cada decisión de diseño
tomada en el capítulo \ref{cap:metodologia}: el dominio de aplicación, las técnicas de
modelado disponibles y las arquitecturas híbridas y de ensamble. De esa línea salen las
respuestas a por qué se agregan los cuatro operadores del CRTM en lugar de una sola
línea, por qué se elige una arquitectura híbrida secuencial y no un ensamble paralelo, y
por qué el protocolo de validación exige un esquema \emph{walk-forward}
\emph{out-of-fold} en lugar de una única partición de retención (\emph{holdout}). La
comparativa final (sección \ref{sec:comparativa_literatura}) frente a nueve trabajos de
referencia delimita el hueco que este TFM se propone ocupar.

\section{Movilidad urbana y demanda de transporte}\label{sec:movilidad_urbana}

La demanda de transporte público en un área metropolitana como Madrid es, por
naturaleza, un fenómeno multimodal: los viajeros se distribuyen entre Metro de Madrid,
la Empresa Municipal de Transportes (EMT), la red de concesionarias de autobús
interurbano (\emph{carretera}) y Renfe Cercanías, y el Consorcio Regional de Transportes
de Madrid (CRTM) es la entidad que agrega y publica esta demanda a escala diaria
para el conjunto del sistema. Sin embargo, la literatura que aborda la predicción de
demanda en Madrid y en ciudades comparables tiende sistemáticamente a operar a una
granularidad muy inferior a la de este TFM: una única línea de autobús
\parencite{monje2022deep}, una única red de metro analizada estación por estación
\parencite{cardozo2012spatial}, o un conjunto multimodal tratado sin la perspectiva de
demanda agregada de todo el sistema \parencite{toque2017short}; incluso cuando se abarca más
de un modo a escala regional, como en el estudio de autobús y ferrocarril de Estambul de
\textcite{yazicioglu2025passenger}, la unidad de análisis sigue siendo la validación horaria
por ruta o región, no el total diario del sistema. \textcite{singhal2014weather}, que
modelan la misma demanda del metro de Nueva York a escala diaria y horaria, encuentran que
ambas escalas difieren sustancialmente en la variabilidad que explican: la granularidad es
una decisión de diseño, no un valor por defecto. Esta elección de escala no es
casual: cuanto más fina es la granularidad —una línea, una estación, una franja
horaria—, más sensible es la serie al ruido idiosincrático de esa unidad concreta (una
obra puntual, un desvío de servicio, un evento singular), mientras que la demanda total
agregada de un sistema de transporte metropolitano completo, al sumar sobre miles de
orígenes y destinos, expone con mucha mayor nitidez los patrones estructurales de fondo
que interesan a un modelo predictivo: la estacionalidad semanal del ciclo
laborable/fin de semana, la estacionalidad anual asociada al calendario académico y
vacacional, y el efecto de los días festivos y ``puente''.

Este TFM adopta
precisamente esa perspectiva sistémica: modela la serie \texttt{total} del CRTM sin
renunciar al seguimiento individual de sus cuatro operadores, que se conserva a lo
largo de todo el pipeline de ingeniería de características (sección
\ref{sec:ingenieria_features}) tanto para el análisis exploratorio como para su uso
como variables de retardo.

Los factores que gobiernan esta demanda agregada, caracterizados en detalle en la
sección \ref{sec:eda}, son de dos naturalezas bien diferenciadas y de peso muy desigual.
Por un lado, la estacionalidad de calendario explica, como respalda la
importancia de variables del capítulo \ref{cap:resultados}, la mayor parte de la
varianza observada: el contraste laborable/sábado/domingo y la posición del día respecto
a un festivo o un puente forman una señal determinista, conocida con años de antelación
y sin incertidumbre de pronóstico. \textcite{yazicioglu2025passenger} replican ese orden
en Estambul: en su clasificación por ganancia de XGBoost encabezan la hora, el día lectivo, el
día de la semana, el festivo y el mes, y la precipitación y la nieve quedan entre las
variables menos relevantes. Por otro lado, las condiciones meteorológicas (temperatura,
precipitación, viento) introducen un efecto de segundo orden, real pero modesto
—\textcite{arana2014weather} lo cuantifican sobre las tarjetas inteligentes de los
autobuses de Gipuzkoa: el viento y la lluvia reducen los viajes y la temperatura los
aumenta—, más difícil de aislar y sujeto además a la incertidumbre propia de cualquier
pronóstico atmosférico, razón por la que este TFM excluye deliberadamente la meteorología
del mismo día del conjunto de variables predictoras (sección \ref{sec:promocion_meteo}).
Entre los trabajos revisados no se ha identificado uno que combine ambos factores,
calendario exhaustivo y física atmosférica retardada de forma anti-fuga, en el marco de
una serie diaria agregada de todo un sistema metropolitano; esa combinación es uno de los
ejes que la sección \ref{sec:comparativa_literatura} formaliza como aporte diferencial.

\section{Modelado de series temporales y aprendizaje automático tabular}\label{sec:modelado_series}

Dos grandes familias de modelos sostienen la arquitectura de este TFM antes de considerar
ninguna hibridación entre ellas. La primera es la de los modelos estadísticos clásicos de series
temporales, fundada por la formulación ARIMA de \textcite{box1970time}: una descomposición autorregresiva, de diferenciación e integración
de medias móviles que, en su extensión estacional SARIMA y con variables exógenas
(SARIMAX), sigue constituyendo siete décadas después una referencia robusta y
difícilmente superable en series con estructura estacional fuerte y bien definida como
la de este proyecto. El SARIMAX$(2,1,2)\times(0,1,2,7)$ seleccionado por rejilla AIC en
la sección \ref{sec:baselines} es, como muestra el capítulo \ref{cap:resultados}, uno
de los dos modelos de una sola etapa más precisos del proyecto, solo superado por los dos
ensambles que los combinan; su fuerza, el modelado explícito de la periodicidad semanal
por diferenciación estacional de orden 7, explica por qué el híbrido no logra desplazarlo.

La segunda familia es la del aprendizaje automático sobre datos tabulares con ingeniería
de características explícita, y en particular los métodos de árboles potenciados por
gradiente, cuyo exponente de referencia actual es XGBoost \parencite{chen2016xgboost}: un
algoritmo de \emph{gradient boosting} regularizado que, a diferencia de una red
neuronal, no requiere aprender representaciones latentes de las variables de entrada y
resulta especialmente eficaz cuando, como en este TFM, la señal más informativa puede
expresarse mediante un conjunto moderado de variables diseñadas a mano (retardos,
medias móviles, términos de Fourier, indicadores de calendario). En demanda de transporte,
\textcite{yazicioglu2025passenger} lo emplean justamente así, como selector de variables por
ganancia previo a las redes recurrentes que comparan. Es precisamente este
algoritmo el que este proyecto emplea en dos papeles distintos: como corrector de
residuos en la etapa 2 de la arquitectura híbrida (sección \ref{sec:etapa2_xgboost}) y,
de forma independiente, como modelo de control entrenado directamente sobre la demanda
total (\texttt{xgboost\_alone}, sección \ref{sec:modelos_control}), cuyo desempeño
resulta ser, junto con SARIMAX, el más sólido entre los modelos de una sola etapa de
todo el estudio.

Sea cual sea la familia de modelo empleada, la validez de su evaluación depende de un
protocolo de partición y validación coherente con la naturaleza secuencial de los datos.
\textcite{bergmeir2012cv} y, más recientemente, \textcite{cerqueira2020evaluating}
muestran empíricamente que los esquemas de validación
cruzada aleatoria estándar, válidos en datos i.i.d., producen estimaciones de error
sistemáticamente optimistas y no comparables entre modelos cuando se aplican a series
temporales, precisamente porque permiten que información del futuro contamine, a través
del reordenamiento de los bloques (\emph{folds}), el ajuste de observaciones pasadas. Este resultado es
la base teórica directa del esquema \emph{walk-forward} expansivo en 5 folds que
genera las predicciones \emph{out-of-fold} de la etapa 1 LSTM (sección
\ref{sec:esquema_oof}) y de la partición estrictamente cronológica 70/15/15 que
gobierna todo el proyecto (sección \ref{sec:particiones_fuga}): ninguna de las dos es una
convención arbitraria, sino la aplicación de un resultado contrastado en la literatura de
evaluación de pronósticos.

Las métricas de
error empleadas en todo el proyecto (MAE, RMSE, MAPE y $R^2$) y los criterios generales
de buenas prácticas en la construcción de baselines siguen, asimismo, el tratamiento de
referencia de \textcite{hyndman2021fpp}. Se reporta el RMSE en
lugar del error cuadrático medio sin raíz por expresar la misma información en las
unidades del propio objetivo y ser, por tanto, directamente comparable con el MAE.

\section{Arquitecturas híbridas y de ensamble}\label{sec:arquitecturas_hibridas}

Los antecedentes de combinar más de un modelo para predecir una misma serie temporal
distinguen dos estrategias que la literatura trata a menudo como intercambiables y que
este TFM mantiene deliberadamente separadas: la hibridación secuencial por corrección de
residuos, y el ensamble paralelo de modelos independientes.

El antecedente directo de la primera estrategia es el trabajo seminal de
\textcite{zhang2003hybrid}, que formula una serie temporal como la suma de una componente
lineal $L_t$ y una componente no lineal residual $N_t$, de modo que
$y_t = L_t + N_t$; ajusta un modelo ARIMA para estimar $\hat L_t$, calcula el residuo
$e_t = y_t - \hat L_t$ y entrena una red neuronal para modelar $e_t$ a partir de sus
propios rezagos, combinando finalmente ambas predicciones como
$\hat y_t = \hat L_t + \hat N_t$. Esta formulación es el precedente
conceptual directo de la arquitectura híbrida evaluada en este TFM
($\hat y^{\text{final}} = \hat y^{\text{LSTM}} + \hat e^{\text{XGBoost}}$, sección
\ref{cap:metodologia}), pero las diferencias de diseño respecto al esquema original de
Zhang son sustanciales y deliberadas, y conviene enunciarlas con precisión. Primero, el
\textbf{orden de las etapas está invertido}: Zhang ajusta primero el componente lineal
(ARIMA) y corrige después con una red neuronal; este TFM ajusta primero el componente no
lineal (LSTM) y corrige después con un modelo de árboles (XGBoost), una inversión
motivada por el hecho de que en este proyecto la componente más difícil de modelar no es
la dinámica lineal de la serie sino precisamente la estructura de calendario, que un
modelo de árboles captura de forma nativa y una red recurrente univariante no puede ver
en absoluto. Segundo, el \textbf{mecanismo de la etapa de corrección es distinto}: en
Zhang, la red neuronal modela el residuo a partir de sus propios rezagos,
$e_t = f(e_{t-1}, \ldots, e_{t-n})$, un esquema univariante; en este TFM, XGBoost modela
el residuo a partir de una matriz de 161 variables predictoras (retardos y estadísticos
móviles de la demanda, rezagos de operador, armónicos de Fourier, calendario y
meteorología retardada) y no de los rezagos del propio residuo, de
modo que la etapa 2 introduce información nueva al sistema y no se limita a extrapolar
la autocorrelación de su propio error.

Tercero, y más relevante para la interpretación
de los resultados del capítulo \ref{cap:resultados}: los residuos que alimentan la
segunda etapa en Zhang son \textbf{residuos \emph{in-sample}} del modelo ARIMA ya
ajustado, mientras que este TFM impone la generación estrictamente
\emph{out-of-fold} de esos residuos mediante el esquema \emph{walk-forward} de 5 folds
descrito en la sección \ref{sec:esquema_oof}. Es una exigencia deliberadamente más
estricta que tiene un coste medible: el 38\,\% de las filas de entrenamiento originales
(889 a 552) se descartan por no disponer de una predicción genuinamente fuera de
muestra. Esta tercera diferencia impide, por diseño, cualquier comparación directa entre
el resultado positivo de Zhang y el negativo de este TFM: son experimentos con un rigor de
validación distinto sobre series de naturaleza distinta, y el capítulo
\ref{cap:resultados} evalúa el suyo bajo el estándar más exigente de los dos.

La hibridación secuencial cuenta también con antecedentes en flujos de pasajeros:
\textcite{chen2019emd} descomponen por modos empíricos (EMD) el flujo entrante de una
estación del metro de Chengdu a intervalos de 15 minutos y alimentan con las componentes una
LSTM que supera a una LSTM directa, a ARIMA y a una red BPN en el día de prueba. Su
protocolo —14 días laborables de entrenamiento, un único viernes de prueba, sin variables
exógenas— ilustra el patrón que el cuadro \ref{tab:comparativa_estado_arte} hace visible:
los híbridos revisados reportan mejoras bajo particiones de un solo corte, sin generar
fuera de muestra las entradas de la segunda etapa.

La segunda estrategia de combinación, el ensamble de modelos independientes que
finalmente resulta ser el mejor modelo de todo este TFM (sección
\ref{sec:por_que_ensemble}), se apoya en un cuerpo teórico distinto: el marco general
de generalización apilada (\emph{stacked generalization}) de
\textcite{wolpert1992stacking}, del que el ensamble de este proyecto es un caso particular
muy simple —un meta-modelo fijo (media aritmética o ponderada) en lugar de uno
aprendido—, y la condición teórica precisa que
determina cuándo esa combinación mejora a sus componentes, formulada por
\textcite{granger1989combining} en su revisión de dos décadas de literatura de combinación
de pronósticos: una combinación solo puede superar a sus partes cuando los modelos
combinados son individualmente subóptimos \emph{y}, de forma crucial, se basan en
conjuntos de información distintos entre sí.

\textcite{bates1969combination}
formalizan además la magnitud exacta de esa mejora: para dos pronósticos con varianza de
error igual $\sigma^2$ y correlación $\rho$ entre sus errores, la varianza del error de
su media aritmética es $\sigma^2(1+\rho)/2$, de modo que cuanto menor es $\rho$, cuanto
menos correlacionados están los errores de los dos modelos combinados, mayor es la
reducción de varianza obtenida, un resultado que
\textcite{perrone1993networks} extienden al contexto de comités de redes neuronales y con
el que este TFM contrasta la ganancia de su ensamble en la sección
\ref{sec:mecanismo_ensemble}.

Esta condición de Granger es exactamente la pieza teórica
que conecta este capítulo con el hallazgo central negativo del capítulo
\ref{cap:resultados} (sección \ref{sec:resultado_negativo_teoria}): la arquitectura
híbrida de este TFM encadena dos etapas cuyos conjuntos de información están anidados
—el pasado reciente de \texttt{total} que ve la etapa 1 está representado en la matriz
de la etapa 2, que además añade el bloque exógeno—, mientras que el
ensamble ganador combina SARIMAX y \texttt{xgboost\_alone}, dos modelos que fallan en
subconjuntos de días sistemáticamente distintos (el primero en laborables ordinarios, el
segundo en fin de semana y festivos) y que, por tanto, sí satisfacen la condición de
información distinta que la hibridación secuencial incumple.
\textcite{makridakis2018m4,makridakis2022m5}, a partir de las competiciones M4 y M5 y de la
constatación reiterada en ambas de que métodos estadísticamente simples y combinaciones de modelos superan con
frecuencia a arquitecturas individuales más complejas (especialmente en muestras de
tamaño moderado como las 1.310 observaciones diarias de este proyecto), proporcionan el
contexto empírico a gran escala que hace de este patrón un resultado esperable y no una
anomalía específica de esta implementación.

\begin{sloppypar}
Por último, el diseño de la matriz de 161 variables predictoras de este TFM —en particular
los rezagos multiescala de demanda y meteorología y los términos armónicos de Fourier
para la estacionalidad semanal y anual (sección \ref{sec:ingenieria_features})— toma
como inspiración metodológica el trabajo de \textcite{surribas_sayago_diffuse}, que
predice radiación solar difusa mediante una fusión paralela CNN+LSTM+MLP con una rama
auxiliar exógena. Conviene fijar el alcance de esa cita: es una referencia
\textbf{metodológica}, de otro dominio de aplicación, y \textbf{no} un antecedente de
predicción de demanda de transporte público. De ella se toman exclusivamente la
construcción de retardos estructurados, la base armónica de Fourier y la idea de una rama
exógena; su arquitectura de fusión paralela no se replica. La corrección secuencial de
residuos de este TFM es un diseño distinto, y así se mantiene tanto en el código como en
esta memoria.
\end{sloppypar}

\section{Análisis comparativo y diferenciación metodológica}\label{sec:comparativa_literatura}

La tabla
\ref{tab:comparativa_estado_arte} sitúa este TFM frente a los nueve trabajos discutidos
en las secciones anteriores sobre cinco ejes comparables: ámbito y granularidad,
horizonte temporal, matriz de variables exógenas, protocolo de validación y aporte
empírico.

\begin{landscape}
  \begin{table}[p]
    \centering
    \footnotesize
    \setlength{\tabcolsep}{4pt}
    \renewcommand{\arraystretch}{0.95}
    \resizebox{\textwidth}{!}{%
      \begin{tabular}{p{2.6cm}p{3.6cm}p{3.0cm}p{4.2cm}p{4.4cm}p{4.6cm}}
\toprule
\rowcolor{naranja}
\textbf{Estudio} & \textbf{(a) Ámbito y granularidad} & \textbf{(b) Horizonte temporal} & \textbf{(c) Matriz exógena} & \textbf{(d) Protocolo de validación} & \textbf{(e) Aporte empírico} \\
\midrule
\rowcolor{filaclara} Monje et al. (2022) & Una línea de autobús (EMT línea 1), solo franja de tarde & 26 meses pre-pandemia (ene. 2015 - feb. 2017) & Festivo binario + calendario (mes, día de semana); cero meteorología & Holdout cronológico simple 80/\allowbreak{}20, una sola ejecución, sin CV & Reporta el ganador (LSTM R²=0,89) sin MAE/\allowbreak{}RMSE/\allowbreak{}MAPE absolutos \\
Toqué et al. (2017) & Transporte multimodal por estación (París) & Logs de billetaje, 2017 & No especificada en la fuente disponible & No especificado en la fuente disponible & Comparativa de métodos de ML para flujos multimodales, sin arquitectura híbrida residual \\
\rowcolor{filaclara} Cardozo et al. (2012) & Entradas por estación del Metro de Madrid, corte transversal espacial & Sin dimensión temporal (regresión espacial) & Entorno construido; sin calendario ni clima & Ajuste espacial in-sample & Regresión espacial por estación; no aborda predicción temporal \\
Zhang (2003) & Series univariantes ajenas al transporte (manchas solares, linces, GBP/\allowbreak{}USD) & 288 /\allowbreak{} 114 /\allowbreak{} 731 observaciones & Ninguna (residuo modelado a partir de sus propios rezagos) & Holdout único, residuos ARIMA in-sample (sin OOF) & El híbrido ARIMA+red neuronal bate a ambos componentes en las tres series \\
\rowcolor{filaclara} Surribas-Sayago et al. [ficha sin verificar] & Radiación solar difusa (dominio meteorológico, no transporte) & No verificado & Retardos y términos de Fourier (fuente de inspiración de la ingeniería de variables de este TFM, no de su arquitectura) & No verificado & Arquitectura paralela CNN+LSTM+MLP; citada solo por su ingeniería de características \\
Este TFM & Demanda diaria agregada de los cuatro operadores del CRTM (Madrid) & 1.310 días continuos 2023-2026, cero huecos, cero nulos, íntegramente post-pandemia & 161 variables: 105 meteorológicas físicas retardadas (lags 1/\allowbreak{}2/\allowbreak{}3/\allowbreak{}7 + media móvil 7), 6 términos de Fourier, 55 días puente detectados; meteorología del mismo día explícitamente prohibida & Partición cronológica 70/\allowbreak{}15/\allowbreak{}15 de frontera única; walk-forward OOF en 5 bloques expansivos; SARIMAX walk-forward de un paso sin reestimación; escalado ajustado solo en train; 257 tests fijando las reglas anti-fuga & Auditoría del fallo: el híbrido no bate a SARIMAX (+70.597 MAE test) ni a XGBoost solo (+78.103); el mejor modelo es un ensamble 50/\allowbreak{}50 (MAE 139.725, -10,5\% vs. el mejor componente) \\
\bottomrule
\end{tabular}
%
    }
    \caption{\headlinecolor{Comparativa metodológica frente a la literatura previa, sobre cinco ejes (ámbito, horizonte, variables exógenas, protocolo de validación, aporte empírico).}}
    \label{tab:comparativa_estado_arte}
  \end{table}
\end{landscape}

De los cuatro primeros ejes, que el cuadro recoge, conviene destacar lo que sitúa a este
trabajo fuera del rango de los estudios de referencia. En \textbf{ámbito} (a) es el único
de los recogidos en el cuadro que agrega los cuatro operadores del CRTM en una serie de
sistema completo, frente a una sola línea de autobús en \textcite{monje2022deep}, a una
estación de metro en \textcite{chen2019emd} y \textcite{cardozo2012spatial}, a rutas y
regiones horarias en \textcite{yazicioglu2025passenger} o a series ajenas al transporte en
\textcite{zhang2003hybrid}. En \textbf{horizonte} (b), sus 1.310 días continuos (2023-2026)
son íntegramente posteriores a la pandemia, a diferencia de los 26 meses prepandemia de
Monje et al., cuya estructura de demanda quedó alterada de forma permanente en 2020-2021.
En \textbf{matriz exógena} (c), ninguno de ellos combina una base meteorológica físicamente
retardada de 105 variables con una codificación de calendario que detecta 55 días puente;
la meteorología del mismo día queda excluida para no asumir un pronóstico atmosférico
perfecto en inferencia. En
\textbf{protocolo de validación} (d) es el único del cuadro que combina partición cronológica de
frontera única, esquema \emph{walk-forward out-of-fold} de 5 folds expansivos para la
etapa 1 y pronóstico SARIMAX \emph{walk-forward} de un paso (un día hacia adelante) sin
reestimación, frente al
\emph{holdout} simple de Monje et al. y Chen et al., los residuos \emph{in-sample} de Zhang o la
regresión OLS explicativa, sin partición temporal, de Singhal et al.; las 257 pruebas
que fijan las reglas anti-fuga del capítulo \ref{cap:datos_features} (de un total de 470,
Anexo \ref{cap:anexo2}) no tienen equivalente declarado en ninguno de los nueve.

El eje decisivo es el \textbf{(e) aporte empírico}, porque es donde este TFM no reporta
una mejora sino la auditoría deliberada de un fallo. A diferencia de Zhang, cuyo híbrido
ARIMA-red neuronal supera a ambos componentes en las tres series que evalúa, y de Chen et
al., cuyo EMD-LSTM bate a sus tres comparadores en un único día de prueba, la
arquitectura híbrida LSTM-XGBoost de este proyecto \textbf{no} bate ni a SARIMAX ni a
\texttt{xgboost\_alone} (capítulo \ref{cap:resultados}), y el modelo finalmente más
preciso es un ensamble de esos mismos dos componentes de una sola etapa (MAE de test
139.725 en su variante equiponderada 50/50 y 139.681 en la ponderada por el inverso del
MAE). Leído a la luz de la condición de Granger (sección
\ref{sec:arquitecturas_hibridas}), ese contraste es en sí mismo la aportación: una
auditoría con validación \emph{out-of-fold} estricta de las condiciones bajo las cuales la
hibridación secuencial \emph{no} mejora sobre sus partes y un ensamble paralelo de modelos
con información parcialmente descorrelacionada sí lo hace.

Entre los trabajos revisados, no se ha identificado un antecedente directo que analice un
resultado híbrido negativo bajo un protocolo equivalente de validación temporal y controles
diagnósticos, sobre una serie de demanda de transporte público real, agregada y multimodal.
Ese es el hueco de investigación (\emph{research gap}) que este TFM ocupa: no la ausencia de arquitecturas híbridas
—existen, Zhang es su fundamento teórico y Chen et al. su aplicación a flujos de
pasajeros—, sino la escasez, en la literatura revisada, de evaluaciones de cuándo y por qué
fallan frente a alternativas más simples.

De ese hueco se derivan las hipótesis de trabajo, enunciadas como las preguntas de
investigación PI1--PI5 de la sección \ref{sec:preguntas_investigacion} con la redacción
exacta con la que se responden en la sección \ref{sec:respuesta_preguntas}, cada una con
un criterio de decisión fijado por escrito antes de observar la métrica correspondiente.
